\documentclass[11pt, oneside]{article} 
\usepackage{geometry}
\geometry{letterpaper} 
\usepackage{graphicx}
	
\usepackage{amssymb}
\usepackage{amsmath}
\usepackage{parskip}
\usepackage{color}
\usepackage{hyperref}

\graphicspath{{/Users/telliott/Dropbox/Github-math/figures/}}
% \begin{center} \includegraphics [scale=0.4] {gauss3.png} \end{center}

\title{Circles}
\date{}

\begin{document}
\maketitle
\Large

%[my-super-duper-separator]

\begin{center} \includegraphics [scale=0.4] {arcs11.png} \end{center}

In the figure above, the angle $s$ is \emph{defined} as the length of arc $a$ that it sweeps out, or subtends, in a unit circle.

In calculus and analytical geometry angles are defined in terms of radians of arc.  For a unit circle with radius = $1$, the total circumference is $2\pi$.

If $s$ is the measure in radians and $D^{\circ}$ the measure in degrees, then
\[ \frac{s}{2 \pi} = \frac{D^{\circ}}{360^{\circ}} \]

It seems natural then to adopt the arc length as a measure of the angle, where $360^\circ$ is equal to $2 \pi$ \emph{radians}, and an angle of $90^\circ$, for example, a right angle, is equal to $\pi/2$ radians.

\begin{center} \includegraphics [scale=0.30] {radian.png} \end{center}

One can easily divide  $360$ by $2 \pi$ to find that one radian is approximately $57^\circ$.
  
To convert some more measures of angles in degrees to radians:
\[ 180^\circ = \pi, \ 90^\circ = \frac{\pi}{2} \]
\[ 60^\circ = \frac{\pi}{3}, \ 45^\circ = \frac{\pi}{4}, \ 30^\circ = \frac{\pi}{6} \]

Central angle and subtended arc are numerically equal, but remember that they are dimensionally different.  Arc is a length, angle is just an angle.

\subsection*{Thales' theorem}

\label{sec:Thales_theorem}

In this chapter, we introduce a few more theorems concerning circles.  We start by revisiting the last of Thales' theorems:

$\bullet$  Any angle inscribed in a semicircle is a right angle.

Think of three points on the circumference of the circle, forming a triangle. If two of the points are on a diameter of the circle, the angle formed at any arbitrary but distinct third point is always a right angle.

$\angle PRQ$ is a right angle.
\begin{center} \includegraphics [scale=0.4] {arcs12.png} \end{center}

\emph{Proof}.

Draw the radius $OR$ (right panel). 

Notice that the two smaller triangles produced ($\triangle OPR$ and $\triangle OQR$) are both isosceles, since two of their sides are radii of the circle.

Therefore, in each triangle the two angles marked with dots of the same color are equal, by the isosceles triangle theorem (\hyperref[sec:isosceles_triangle_theorem]{\textbf{ref}}).

Since $\angle PRQ$ contains one angle of each type, it is equal to one-half the angle sum for the triangle, i.e., $\angle PRQ$ is a right angle.

To restate this in more conventional notation:  in the figure below, $\angle PRQ = \phi + \theta$.  

\begin{center} \includegraphics [scale=0.4] {arcs13.png} \end{center}

Since the full measure of the triangle is $180^\circ = \pi$ radians, and
\[ \phi + \phi + \theta + \theta = \pi \]
it follows that
\[ \phi + \theta = \pi/2 \]

$\square$

\subsection*{angles on the perimeter}

\label{sec:peripheral_angle}

The arc swept out by the angle $\angle POQ$ (equal to two right angles) is clearly equal to $\pi$, because that arc is one-half of a circle.  

But the two angles on the perimeter of the circle which together subtend one-half of the circle arc add up to
\[ \angle PRQ = \phi + \theta = \pi/2 \]

What's going on?

To clarify, let us label the angles at the center of the circle, as well as three segments of arc.

\begin{center} \includegraphics [scale=0.4] {arcs8.png} \end{center}

$a$ is the arc swept out by angle $s$, and $a$ and $s$ have the same measure by definition, although one is a length and the other an angle.

\[ s = a \]

By the external angle theorem (\hyperref[sec:external_angle_theorem]{\textbf{ref}}), we know that
\[ s = 2 \phi \]
and so conclude that
\[ \phi = \frac{s}{2} \]

The angle $\phi$ lying on the perimeter sweeps out the same arc as $s$, even though $\phi$ is one-half of $s$.

$\square$

$\phi$ is farther away from the arc, so we get a bigger arc for the same angular measure.  Another way to say the same thing:  $\phi$ is farther away from the arc, so we get a smaller angle for the same arc.

$\bullet$  An angle on the perimeter (what we will call a peripheral angle) is twice the central angle subtending the same arc.  This is called the \emph{inscribed angle theorem}

A related approach is to simply count up the angles in two triangles.  For the large right triangle we have

\[ \pi = \phi + \phi + \theta + \theta \]

and for the smaller triangle with central angle $s$

\[ \pi = s + \theta + \theta \]

Clearly, then, $s = 2 \phi$ so

\[ \phi = \frac{s}{2} \]

$\square$

Here is a different proof of the inscribed angle theorem for the case where the peripheral angle does not include a diameter of the circle.  It uses various isosceles triangles and the sum of angles theorem.

By definition, the central angle subtending a given arc is equal in measure to the length of the arc.  We claim that the peripheral angle subtending the same arc is one-half the central angle so  $2 \phi = \theta$.

\begin{center} \includegraphics [scale=0.5] {broken_chord1.png} \end{center}

\emph{Proof}.

The figure has two isosceles triangles, where the equal sides are all radii of the circle.  By the triangle sum theorem and the isosceles triangle theorem:

\[ 2(A + \phi) + C = \pi \]
\[ 2A + \theta + C = \pi \]

Equating the two results and canceling we get
\[ 2A + \theta = 2(A + \phi) \]
\[ \theta = 2 \phi \]

$\square$

A proof using this approach for the case where the peripheral angle includes the center of the circle is left as an exercise.  It is extremely easy.

\subsection*{theorem on cyclic quadrilaterals}

\label{sec:quadrilateral_supplementary}

All this leads to a wonderful simple theorem about quadrilaterals.

\begin{center} \includegraphics [scale=0.4] {circles_4.png} \end{center}

$\bullet$ \ For \emph{any} quadrilateral whose four vertices lie on a circle, the opposing angles are supplementary (they sum to $180^\circ$).

\emph{Proof}.

Together, opposing angles exactly subtend the whole arc of the circle.

$\square$

\subsection*{tangent}

$\bullet$  The arc swept out between a tangent and a chord is equal to the arc lying between the point of tangency and the point where the arc meets the circle.

Acheson calls this the alternate segment theorem.

\begin{center} \includegraphics [scale=0.4] {arcs14.png} \end{center}

\emph{Proof}.

The arc $a$ is equal to twice the arc swept out by angle $\theta$ on the far right of the diameter.  But the angle between the tangent and the chord is also equal to $\theta$, because $\theta + \phi$ is equal to a right angle.

Therefore the arc swept out between the tangent and the chord is equal to $a$.

$\square$

\subsection*{geometric mean}

We showed in the chapter on the Pythagorean theorem (\hyperref[sec:pythagorean_thm]{\textbf{ref}}) that the altitude of a right triangle is the geometric mean of the two components of the base.

\[ h^2 = pq \]
\[ h = \sqrt{pq} \]

According to wikipedia:

\url{https://en.wikipedia.org/wiki/Geometric_mean}

The fundamental property of the geometric mean is that (letting $m$ be the \emph{geometric mean} here):
\[ m \ [ \ \frac{x_i}{y_i} \ ] \ = \frac{m(x_i)}{m(y_i)} \]

and one consequence is that

\begin{quote}This makes the geometric mean the only correct mean when averaging normalized results; that is, results that are presented as ratios to reference values.\end{quote}

A number of examples are given in the article.

This section is here because originally, there was a proof-without-words that the geometric mean is always less than or equal to the arithmetic mean.

\begin{center} \includegraphics [scale=0.4] {arcs15.png} \end{center}

I decided to add some words.  A right triangle is inscribed in a semicircle.  

Let the length of the long dotted line be $c$.  Then
\[ a^2 + h^2 = c^2 \]
Let the length of the short dotted line be $d$.  Then
\[ b^2 + h^2 = d^2 \]

Adding together
\[ a^2 + b^2 + 2 h^2 = c^2 + d^2 \]

But a third application of the Pythagorean theorem shows that the right-hand side is just $(a + b)^2$ so
\[ a^2 + b^2 + 2 h^2 = (a + b)^2 \]
\[ 2 h^2 = 2 ab \]
\[ h^2 = ab \]

The square of the altitude $h$ is equal to the product of chord segments (we will prove this geometrically in the next chapter as well).

\[ h = \sqrt{ab} \]

But we also have that $a + b = 2r$ and hence
\[ r = \frac{a + b}{2} \]
Do you recognize these?  The second expression is the arithmetic mean of $a$ and $b$, while the first is the geometric mean.

The geometry shows that $h \le r$ so:
\[ \sqrt{ab} \le \frac{a + b}{2} \]

The geometric mean is always less than the arithmetic mean, except when $a = b$, where they are equal (or all of $n$ values are equal).

\subsection*{problem}

\begin{center} \includegraphics [scale=0.3] {circles1.png} \end{center}

Two circles meet at $Q$ and $S$.  $QR$ and $QT$ are diameters of the two circles.  Prove that $RST$ are colinear.

\begin{center} \includegraphics [scale=0.3] {circles2.png} \end{center}

Since $QR$ is a diameter of the circle centered at $O$, $\angle QSR$ is a right angle.  

And so is $\angle QST$, since $QT$ is a diameter of the second circle.  

Hence the total angle at $S$ is two right angles or a straight line.  Therefore $RS$ and $ST$ togethere form a straight line segment.

\subsection*{double arc problem}

This problem is taken from an online collection by David Surowski.

\url{https://www.math.ksu.edu/~dbski/writings/further.pdf}

\begin{center} \includegraphics [scale=0.6] {further1.png} \end{center}

Given that $AB$ and $AC$ are tangents to the circle meeting at $A$.  Given a second tangent $DE$, meeting the circle at $P$.  Prove that the arc subtended by $\angle BOC$ is twice that subtended by $\angle DOE$.

\emph{Proof}.

Notice that $DB$ and $DP$ are tangents to the circle meeting at $D$.  Therefore $DB = DP$ and then $\triangle BOD \cong \triangle DOP$, so $\angle BOD = \angle DOP$.

The same argument applies to $EC$ and $EP$.  Therefore the inner arc is composed of two angles, while the outer arc has two copies of each of those angles.

$\square$

\subsection*{problem}

\begin{center} \includegraphics [scale=0.4] {Posamentier1_4.png} \end{center}

Posamentier gives this problem.  I is the center of the circle (the circumcircle of $\triangle ABC$).  $AE$ is the altitude of $\triangle ABC$.  

\emph{To prove}.

If the angle at vertex $A$ is bisected (by $AD$), then $AD$ also bisects $\angle EAI$.

First of all, I would restate the problem:  prove that $\angle BAE = \angle CAI$.  Then, the bisector of $A$ must bisect $\angle EAI$.

Hint:  draw $IC$ and use the relationships from this chapter as well as complementary angles in a right triangle.

\subsection*{problem}

Here is another problem from Surowski.  This one is fairly easy, with one assumption, which I'm pretty sure is justified.

Given a circle with center at $O$.  Points $XAFC$ are colinear and tangent at $A$.  Points $EOBDC$ colinear as well.  $FD$ is tangent at $D$.

\begin{center} \includegraphics [scale=0.35] {Surowski_p30_5.png} \end{center}

To prove:  $AE$ bisects $\angle BAX$ and $AD$ bisects $\angle BAC$.

\emph{Proof}.

We are not given that $AB \perp ED$.  On the other hand, if $AB$ is not constrained, I don't know how to solve the problem, and if $AB$ is not $\perp ED$, then the angle bisection we're supposed to prove would not be correct.

$ED$ goes through the center $O$, so it is a diameter of the circle.  $\triangle EAD$ is inscribed in a semicircle, hence $\angle EAD$ is a right angle.

\begin{center} \includegraphics [scale=0.5] {Surowski_p30_5a.png} \end{center}

$\angle AEB$ and $\angle ADE$ are complementary.  Since $FD \perp EDC$, $\angle ADF$ and $\angle ADE$ are complementary, so $\angle ADF = \angle AEB$. 

Since $FA$ and $FD$ are both tangents from the point $F$, they are equal so $\triangle ADF$ is isosceles, and $\angle FAD = \angle ADF$.

To make further progress, we must assume $AB \perp ED$.

\begin{center} \includegraphics [scale=0.5] {Surowski_p30_5b.png} \end{center}

Then, the two component right triangles are similar to $\triangle ADE$, so the angles at vertex $A$ of $\triangle AED$ are complementary and equal as marked. 

Finally, $\angle XAE$ is equal to $\angle EAB$ because two pairs of complementary angles add up to two right angles.

$\square$

\subsection*{Queen Dido}

The mighty city of Carthage was the capital city of the Phoenicians.  As Rome grew strong, there was a titanic struggle between the two peoples, which Carthage eventually lost.  The ruins of Carthage lie near present-day Tunis.

Queen Dido was the legendary founder of the the city of Carthage.  She was supposedly 

\begin{quote}granted as much land as she could encompass with an oxhide.  She promptly cut the ox-hide into very thin strips.\end{quote}

The problem then is to maximize the area enclosed by a curve of fixed length.

\begin{center} \includegraphics [scale=0.5] {Dido.png} \end{center}

In calculus there are a number of problems like this.  What's nice is that this problem has a wonderful solution that uses only the tools we have so far.  In particular, we need the \hyperref[sec:Thales_circle_theorem_converse]{\textbf{converse of Thales circle theorem}}.

The argument goes as follows.  Suppose we have a particular outline for the city limits and we're pretty happy with it.  We suppose it is a maximum (left panel).

\begin{center} \includegraphics [scale=0.5] {Dido2.png} \end{center}

Then we notice that by rearranging $AD$ and $BD$ so they meet at a right angle, the crescent-shaped areas are unchanged, but the area of $\triangle ABD$ is a maximum.  That's because a right triangle, having the two sides at right angles, has area equal to the product of the two sides (divided by $2$).  No other triangle with the same two sides has as much area.

So the arrangement on the right has a bigger area.

But then, with $AB$ as the diameter of a circle, if $\angle ADB$ is a right angle, it must lie on the circumference of that circle, by the converse of Thales' theorem.

And this is true regardless of the relative lengths of $AD$ and $BD$.  Therefore the maximum area is obtained when $D$ traces out a semi-circle.

This example is in Acheson's Geometry.


\end{document}